\documentclass[../ap_2014.tex]{subfiles}

\graphicspath{{./image/}}

\begin{document}

\setcounter{chapter}{2}
\chapter{物性:フォノン}
\section{}
系の次元を\(d\), フォノン分散を\(\omega_k\)とする。
状態密度は
\begin{equation}\begin{aligned}[b]
    D(\omega)
    &=\sum_k\delta(\omega-\omega_k)\\
    &\propto \int \dd[d]{k} \delta(\omega -vk)\\
    &\propto \int \dd{k} k^{d-1}\delta\qty(k-\frac{\omega}{v})\\
    &\propto \omega^{d-1}
\end{aligned}\end{equation}
これより
(a)3次元立方格子系では\(p=2\)で、
(b)2次元立方格子系では\(p=1\)である。

\section{}
フォノンはボーズ分布に従うので、
低温での内部エネルギーと比熱は
\begin{equation}\begin{aligned}[b]
    E &= \int_{0}^{\infty}\frac{\hbar\omega}{e^{\hbar\omega/k_BT}-1}D(\omega)\dd{\omega}\\
    &\propto (k_BT)^{d+1}\int_{0}^{\infty}\\
    C &= \dv{E}{t} = T^d
\end{aligned}\end{equation}
とわかる。
これより
(a)3次元立方格子系では\(q=3\)で、
(b)2次元立方格子系では\(q=2\)である。

\section{}
古典的なバネの描像を使うと、バネ定数を\(k\)とすると固有振動数は\(\omega\propto\sqrt{k}\)となる。
層間間の結合は弱いため、\(\omega\)は小さくなるので、
(i)が層間、(ii)が層内であると分かる。

\section{}
\(\omega_0\ll \omega_a\)のときでは、
\(k'\)方向が正方形で\(k\)方向が長方形の直方体の表面である。

\(\omega_a\ll \omega_0 \ll \omega_b\)のときでは、
\(k'\)方向が正方形で\(k\)方向が長方形の直方体の表面のうち、正方形の面を取り除いたものである。

\section{}
\(\omega_0\ll \omega_a\)のときでは、
全ての方向のモードが励起するので、
状態密度は3次元のときの\(D(\omega)=\omega^2\)となる。
なので、\(r=2\)である。

\(\omega_a\ll \omega_0 \ll \omega_b\)のときでは、
層間のモードはなく、層内の正方形のモードだけが励起するので、
状態密度は2次元のときの\(D(\omega)\omega\)となる。
なので、\(r=1\)である。

\section{}
\(\omega_0\ll k_BT\)のときでは、
全ての方向のモードが励起するので、
状態密度は3次元のときの比熱\(C\propto T^3\)となる。
なので、\(s=3\)である。

\(\omega_a\ll k_BT \ll \omega_b\)のときでは、
層内の方向のモードのみが励起するので、
状態密度は2次元のときの比熱\(C\propto T^2\)となる。
なので、\(s=2\)である。

\(\omega_b \ll k_BT\)のときでは、
フォノン励起は見えなくなる。
このときには等エネルギー分配則が成り立つので、
内部エネルギーは\(T\)に比例する。
なので比熱は定数となり、\(s=0\)である。


\section*{感想}
冪しか気にしてないあたり、理論の人の作った問題のような気がする。
フォノンと次元の関係を整理できてよかった。


\end{document}
